%% icmi2026_label_augmentation.tex
%% Target: 8-page ICMI/CSMI workshop paper.
%%
%% Compile:
%%   pdflatex icmi2026_label_augmentation
%%   bibtex   icmi2026_label_augmentation
%%   pdflatex icmi2026_label_augmentation
%%   pdflatex icmi2026_label_augmentation

\documentclass[sigconf,review,anonymous]{acmart}

\usepackage{booktabs}
\usepackage{multirow}
\usepackage{amsmath}
\usepackage{enumitem}
\usepackage{array}
\usepackage{xcolor}
\usepackage{microtype}

% Layout tuning: allow inter-word stretch of up to 3em before declaring
% overfull (standard practice for citation-dense conference papers).
\setlength\emergencystretch{3em}
% Suppress underfull hbox warnings for short paragraph-ending lines.
\hbadness=10000

\acmConference[CSMI '26]{Workshop on Collective States in Multimodal Interaction (CSMI)}
  {October 2026}{Naples, Italy}
\acmYear{2026}
\copyrightyear{2026}
\settopmatter{printacmref=true}
\pagestyle{plain}

\definecolor{bestcol}{HTML}{1a6e1a}
\newcommand{\best}[1]{\textbf{\textcolor{bestcol}{#1}}}

\title{Personality-Weighted Label Augmentation for Physiological Affect Sensing
in Small Collaborative Groups}

\author{Anonymous Authors}
\affiliation{\institution{Anonymous Institution}}
\email{anon@example.com}

\begin{document}

\begin{abstract}
Physiological affect sensing in naturalistic group interaction is limited less
by sensor density than by label density: wearable devices produce thousands of
time windows, while self-report labels are collected only a few times per
session.
This paper asks whether stable personality similarity can serve as a trust
signal for pseudo-label augmentation in small collaborative groups.
Using GroupAffect-4, a four-person collaborative dataset with wearable
physiology, eye tracking, Big Five personality, and post-task VAD labels, we
compare Gaussian-process, physiology-derived, and personality-mediated
augmentation variants under a shared target-construction pipeline.
An RBF-SVM with balanced midpoint-tertile labels reaches $0.495$ macro-F1 on
the canonical known-team temporal holdout.
Replacing GP confidence with BFI-44 cosine similarity as the pseudo-label
trust weight improves macro-F1 to $0.526$, with the largest gains in Valence
($+0.038$) and Dominance ($+0.050$).
In a cross-subject LOSO stress test (no augmentation), aggregate macro-F1 is
$0.320$ (V$=0.190$, A$=0.438$, D$=0.334$), showing a substantial transfer gap
relative to the known-team protocol.
However, strict session-isolated LOGO evaluation removes the augmentation
benefit ($0.262$ vs. $0.266$ no augmentation), showing that personality
weighting is a within-team augmentation mechanism rather than an unseen-group
transfer solution.
Across held-out sessions, AP1/AP2 differences are not significant against A0
(paired Wilcoxon signed-rank, all $p>.2$).
Exploratory post-block questionnaire correlations provide external validity:
predicted arousal tracks engagement ($r=0.46$) and predicted valence tracks
satisfaction ($r=0.39$).
\end{abstract}

\ccsdesc[500]{Computing methodologies~Supervised learning by classification}
\ccsdesc[300]{Human-centered computing~Collaborative interaction}
\ccsdesc[200]{Applied computing~Life and medical sciences}

\keywords{label augmentation, physiological affect sensing, Big Five
personality, VAD prediction, Gaussian Process, group interaction}

\maketitle

\section{Introduction}
\label{sec:intro}

Wearable physiological sensors provide dense streams of electrodermal,
cardiac, oculomotor, pupil, and movement signals, but affect labels remain
sparse in naturalistic interaction.
In collaborative group studies, participants can only be interrupted for
self-report a few times per session.
The result is a structural mismatch: thousands of physiological windows, but
only a small number of labelled Valence--Arousal--Dominance (VAD) observations.

The usual response is to interpolate labels over time or to train larger
models on small labelled sets.
Both strategies are fragile.
Physiological affect is strongly person-specific, and the same arousal pattern
can reflect different appraisals depending on personality, task role, and group
context.
We therefore test a simple hypothesis: in small collaborative groups, stable
personality similarity may be a better trust signal for cross-person
pseudo-labels than instantaneous physiological synchrony.

This paper focuses on one question: when does personality-weighted
pseudo-labelling help sparse-label physiological affect modelling, and when
does it fail under stricter transfer protocols.
The paper is intentionally scoped around low-capacity supervised models.
We use the same RBF-SVM throughout the main comparison so that improvements can
be attributed to pseudo-label construction rather than to architecture changes.
This is important because the dataset has only 103 labelled training windows in
the canonical split; high-capacity models can easily turn an augmentation
comparison into a capacity comparison.

The practical motivation is a known-team monitoring setting.
A team completes multiple collaborative tasks, physiology is recorded
continuously, and only sparse self-report is available.
The system may use earlier-task data from that same team to support later-task
prediction.
This is not the same as predicting a wholly unseen group, and much of the
paper is devoted to making that boundary explicit.
Our contributions are:

\begin{itemize}[leftmargin=1.2em,itemsep=1pt,topsep=2pt]
  \item We define a leakage-aware augmentation taxonomy for sparse VAD labels
  in small groups, separating known-team temporal holdout from LOSO
  cross-subject and strict LOGO unseen-group stress tests.
  \item We compare GP, EDA-derived, kNN-derived, and personality-mediated
  augmentation variants under the same RBF-SVM base model.
  \item We show that BFI-44 cosine similarity is the strongest trust weight in
  the known-team setting ($0.526$ macro-F1), but does not improve strict LOGO.
  \item We validate the resulting predictions against independent post-block
  questionnaire items not used for training.
\end{itemize}

\begin{sloppypar}
\section{Related Work}
\label{sec:related}

\paragraph{Dimensional affect and target construction.}
Russell's circumplex model foregrounds Valence and Arousal
\cite{russell1980circumplex}, while the PAD tradition adds Dominance as a
third interpersonal dimension \cite{mehrabian1996pleasure}.
The Self-Assessment Manikin (SAM) operationalises these axes as compact
visual self-reports \cite{bradley1994sam}, but ratings are ordinal, tied,
and often task-dependent.
Fixed threshold binning produces class imbalance that biases classifiers
\cite{he2009imbalanced,johnson2019classimbalance}; label distribution and
ordinal learning approaches argue for preserving ordered uncertainty
\cite{geng2016label,wen2023ordinal}, motivating our balanced
midpoint-tertile target pipeline.

\paragraph{Physiological and oculomotor affect recognition.}
DEAP~\cite{koelstra2012deap}, DREAMER~\cite{katsigiannis2018dreamer},
MAHNOB-HCI~\cite{soleymani2012mahnob}, WESAD~\cite{schmidt2018wesad},
and ASCERTAIN~\cite{subramanian2018ascertain} established multimodal
benchmarks; reviews confirm EDA, cardiac, and wearable signals are
useful but sensitive to preprocessing and splitting
protocol~\cite{shu2018review,ahmad2022survey,kreibig2010autonomic}.
PPG and EDA processing follow NeuroKit2~\cite{makowski2021neurokit2} and
cvxEDA~\cite{greco2016cvxeda} respectively.

\paragraph{Social and group affect datasets.}
IEMOCAP~\cite{busso2008iemocap} and RECOLA~\cite{ringeval2013introducing}
introduced acted and remote dyadic affective interaction;
K-EmoCon~\cite{park2020kemocon} added naturalistic dyadic debate with
wearable sensors and multiple annotation perspectives.
AMIGOS~\cite{mirandacorrea2021amigos} explicitly studies affect, personality,
mood, and social context, but the group condition is shared media viewing
rather than collaborative problem solving.
Meeting and social-signal corpora such as AMI~\cite{mccowan2005ami},
SALSA~\cite{alameda2015salsa}, MatchNMingle~\cite{cabrera2018matchnmingle},
and ELEA~\cite{sanchez2012emergent} capture rich group behaviour and
social signals.
GroupAffect-4~\cite{groupaffect4} occupies the remaining intersection:
structured four-person collaboration with per-person physiology, eye
tracking, VAD self-report, and pre-session personality.

\paragraph{Sparse labels, augmentation, and generalisation.}
Small labelled physiological datasets make model selection unusually
dependent on target quality and validation design.
Time-series augmentation can improve wearable classifiers
\cite{um2017data,iwana2021augmentation}, and contrastive methods such as
TS-TCC~\cite{eldele2021tscc} learn representations from unlabelled
segments.
Semi-supervised consistency methods, from Mean
Teacher~\cite{tarvainen2017mean} to FixMatch~\cite{sohn2020fixmatch},
show that pseudo-labels are useful only when confidence is reliable;
pseudo-labelling adds unlabelled examples above a confidence threshold
\cite{lee2013pseudo}.
Gaussian Processes provide uncertainty-aware interpolation in label
space~\cite{rasmussen2006gaussian,atcheson2017gp}, complementing
confidence-threshold approaches.
Subject- and group-aware splits are essential for interpreting wearable
affect model generalisation~\cite{ahmad2022survey,ventura2022validation}.

\paragraph{Personality and affect dynamics.}
The Big Five~\cite{john1999bigfive} is the dominant personality taxonomy;
Neuroticism and Extraversion are primary affect predictors
\cite{watson1992traits}.
ASCERTAIN and AMIGOS show that personality shapes physiological affect
responses~\cite{subramanian2018ascertain,mirandacorrea2021amigos}.
\citet{meng2026moderating} used CTSEM to model how all five traits moderate
VAD dynamics during real interaction, finding trait-specific moderation
timescales that directly motivate our per-dimension personality trust
weighting while also warning that trait effects are contextual and slow
relative to 15-second windows.
Emotional contagion suggests group members converge in affective state
\cite{hatfield1993emotional}, supporting cross-person pseudo-label transfer
weighted by personality similarity.

\paragraph{Why personality and not only physiological synchrony?}
Physiological synchrony is measured at the same timescale as the target
windows, but it is ambiguous: co-fluctuation can reflect a shared external
event, concurrent movement, or genuine affective convergence.
Personality similarity is slower and more stable, providing a durable prior
for weighting cross-person pseudo-labels when self-reports are sparse and
uncertain.
Our contribution is to operationalise BFI cosine similarity as a trust
weight and to test its deployment boundary under known-team, LOSO, and
strict LOGO protocols.
\end{sloppypar}

\section{Dataset and Targets}
\label{sec:dataset}

GroupAffect-4 contains 10 four-person collaborative sessions, with 39
participants retained after data completeness checks \cite{groupaffect4}.
Each session contains five sequential tasks: T0 passive rest, T1 information
pooling, T2 negotiation, T3 ideation, and T4 public-goods interaction.
After each task, participants reported Valence, Arousal, and Dominance using
9-point SAM scales \cite{bradley1994sam}; T4 Dominance was not administered
by design.
Participants also completed the BFI-44 Big Five Inventory
\cite{john1999bigfive} before the session.

\paragraph{Physiology and eye-tracking inputs.}
Table~\ref{tab:modalities} summarises the two wearable hardware platforms.
All streams are time-stamped via LSL (Lab Streaming Layer) for cross-modal
alignment and yield a 49-dimensional summary feature vector per 15-second
window (mean, standard deviation, slope, and percentile statistics per
channel).
PPG is processed with NeuroKit2~\cite{makowski2021neurokit2}; EDA is
decomposed into phasic (SCR) and tonic (SCL) components via
cvxEDA~\cite{greco2016cvxeda}; gaze/pupil statistics follow standard
fixation and pupillometry conventions~\cite{laeng2012pupillometry}.

\begin{table}[t]
\caption{Physiology and eye-tracking inputs (GroupAffect-4).}
\label{tab:modalities}
\centering
\setlength{\tabcolsep}{3pt}
\footnotesize
\begin{tabular}{>{}l >{}l >{}l >{}l}
\toprule
\textbf{Modality} & \textbf{Device} & \textbf{Rate} & \textbf{Features} \\
\midrule
PPG        & EmotiBit & $\sim$25\,Hz & HR, RMSSD, SDNN, amp. \\
EDA        & EmotiBit & $\sim$25\,Hz & SCL/SCR, phasic, tonic \\
Skin temp. & EmotiBit & $\sim$25\,Hz & Mean, slope \\
IMU        & EmotiBit & $\sim$25\,Hz & Accel.~mag., movement \\
Gaze       & Tobii G3 & 40--60\,Hz   & Fix.~rate, saccade \\
Pupil      & Tobii G3 & 40--60\,Hz   & Mean diam., dilation var. \\
\bottomrule
\end{tabular}
\end{table}

We use 15-second windows with 50\% overlap.
The canonical split is a known-team temporal holdout:
train on T0+T1 ($N=103$ windows), validate on T2 ($N=78$), and test on T3
($N=65$).
This design keeps all participants represented in every split while preserving
task order.
Sliding 15-second windows with 50\% overlap yield 292 labelled windows and
8{,}221 unlabelled windows across between-task intervals.
We also report LOSO and LOGO stress tests for cross-person and cross-group
generalisation.
All main results use macro-F1 per VAD dimension and report the three
dimensions separately.
The mean score is useful for ranking variants, but it can hide the main
scientific pattern: Valence and Dominance benefit most from personality
weighting, while Arousal depends more on physiological timescale and temporal
proximity to self-report.

\paragraph{Target construction.}
Raw SAM scores are converted into Low, Mid, and High classes using midpoint
tertile thresholds fitted on training data only.
Boundaries are set at midpoints between integer response values.
This keeps tied Likert clusters in the same class.
Per-dimension inverse-frequency weights correct residual imbalance.
This target pipeline replaces a fixed-threshold baseline that inflated
Valence F1 when the T3 test set contained no Low-valence examples.
The same fitted thresholds are used to discretise GP posterior moments for
augmented windows, so hard labels and pseudo-labels live in the same class
space.

\subsection{Why Target Construction Matters}
\label{sec:target_construction}

The label pipeline is part of the contribution because augmentation quality
cannot be evaluated independently of target quality.
The original 1--9 SAM responses are ordinal, tied, and task-dependent.
For example, T3 ideation contains many positive-valence reports; fixed
thresholds such as Low $\leq 3$, Mid $=4$--$6$, and High $\geq 7$ can remove
one class from the test fold.
When that happens, macro-F1 is no longer a three-class metric even if the code
reports a single number.

Balanced midpoint-tertile thresholds avoid this failure mode.
The split points are fitted on training labels only, then moved to midpoints
between integer response values so that all tied Likert responses stay in the
same class.
This gives a small but honest benchmark: the SVM baseline becomes stronger
because class semantics match the task distribution, while the MLP no longer
benefits from silent class collapse.
All augmentation variants in this paper are evaluated after this correction.

\begin{table}[t]
\caption{Label-pipeline ablation motivating the shared target construction.
Values are T3 macro-F1 for an MLP without augmentation.}
\label{tab:target_ablation}
\centering
\setlength{\tabcolsep}{4pt}
\small
\begin{tabular}{lccc}
\toprule
\textbf{Label pipeline} & \textbf{V} & \textbf{A} & \textbf{D} \\
\midrule
Fixed thresholds, no weights & 0.575 & 0.516 & 0.338 \\
Balanced tertile, no weights & 0.372 & 0.443 & 0.429 \\
Full target pipeline         & 0.405 & 0.464 & 0.495 \\
\bottomrule
\end{tabular}
\end{table}

\begin{table}[t]
\caption{Canonical benchmark setup used throughout this paper.}
\label{tab:setup}
\centering
\setlength{\tabcolsep}{4pt}
\small
\begin{tabular}{ll}
\toprule
\textbf{Component} & \textbf{Definition} \\
\midrule
Dataset & GroupAffect-4, 39 participants, 10 groups \\
Signals & Gaze, pupil, EDA, PPG, skin temp., IMU \\
Features & 49 physiology/eye-tracking summary features \\
Labels & Balanced midpoint-tertile VAD classes \\
Train/val/test & T0+T1 / T2 / T3 \\
Base model & RBF-SVM, $C=1$, balanced class weights \\
Metric & Macro-F1 per VAD dimension and mean \\
\bottomrule
\end{tabular}
\end{table}

\subsection{Deployment Regimes}
\label{sec:deployment}

We distinguish three regimes because ``generalisation'' has different meanings
in group affect sensing.
\textbf{Known-team temporal holdout} asks whether earlier observations from a
team can help predict that same team's later task windows.
This is realistic for longitudinal team-monitoring systems and is the primary
setting for AP1.
\textbf{LOSO} asks whether models transfer to an unseen participant within the
same task family.
\textbf{LOGO} asks whether models transfer to an entirely unseen group, where
all members, baselines, and interaction dynamics are absent from training.
The last regime is much harder and should not be conflated with the first.

This separation is important for personality-based augmentation.
BFI-44 is collected before the session, so the personality similarity weight is
not label leakage.
However, a personality weight applied to a pseudo-label does not make the
pseudo-label itself automatically leakage-free.
The strict LOGO experiment therefore regenerates the pool while excluding the
held-out group entirely.

\section{Augmentation Methods}
\label{sec:methods}

All augmentation variants use the same SVM base model.
Unlabelled windows are pseudo-labelled and assigned sample weights; test
evaluation always uses ground-truth SAM-derived labels only.
Let $\mathcal{D}_{L}=\{(\mathbf{x}_i,y_i,w_i=1)\}$ denote labelled training
windows and $\mathcal{D}_{U}=\{(\mathbf{x}_u,\hat{y}_u,\alpha_u)\}$ denote
accepted pseudo-labelled windows.
The SVM is trained on
\begin{equation}
  \mathcal{D}_{aug} = \mathcal{D}_{L} \cup \mathcal{D}_{U},
\end{equation}
with sample weight $\alpha_u \in [0,1]$ for pseudo-labelled examples.
For AP1, $\alpha_u$ is the BFI similarity weight; for AP2 it is the product of
BFI similarity and GP posterior confidence.
No test labels are used in pseudo-label generation.

\paragraph{GP pseudo-labels.}
For each participant and VAD dimension, we fit an Ornstein--Uhlenbeck GP to
task-level SAM observations.
The posterior over the 1--9 SAM scale is integrated across the
training-fitted midpoint thresholds to produce a three-class soft label.
For the SVM, the pseudo-label class is the posterior argmax and the
posterior confidence becomes a sample weight.
Threshold alignment is critical: the GP posterior is first expressed on the
original 1--9 SAM scale and only then integrated over the training-fitted
Low/Mid/High thresholds.
Using stale thresholds from a different target pipeline can invert the effect
of augmentation because the pseudo-label class prior no longer matches the
supervised labels.

\paragraph{Physiology-derived arousal labels.}
Because Arousal GP calibration is weak, we test two alternatives:
EDA Ridge regression and kNN labels in the full 49-dimensional feature space.
The EDA-only approach regresses toward the mean; kNN improves leave-one-out
arousal accuracy but still introduces noisy decision-boundary samples.

\paragraph{Personality trust weights.}
For participant $i$ in group $G$, we compute BFI cosine similarity to the
other group members:
\begin{equation}
  \mathrm{BFI\_sim}_i =
  \frac{1}{|G|-1}\sum_{j \neq i}
  \frac{\mathbf{b}_i \cdot \mathbf{b}_j}{\|\mathbf{b}_i\|\|\mathbf{b}_j\|}.
\end{equation}
AP1 uses $\mathrm{BFI\_sim}_i$ as the pseudo-label trust weight.
AP2 multiplies this personality weight by GP posterior confidence.
The BFI weight itself is leakage-free because personality is collected before
the session, but the pseudo-label class inherits whatever information was used
to generate the GP pool.

\paragraph{Leakage taxonomy.}
Scenario 1 is the known-team temporal holdout: the pool is filtered to
training tasks, but pre-generated GP parameters may encode mild distributional
information from the full dataset.
Scenario 2 is session-isolated LOGO: the held-out session is excluded from
pool generation, BFI maps, and parameter estimation.
Only Scenario 2 tests unseen-group deployment.

\paragraph{Pool filtering.}
Pool windows are filtered by task and source self-report index.
For the known-team temporal split, T3 labels are excluded from augmentation
because T3 is the test task.
For LOGO, all windows from the held-out group are removed before GP fitting,
BFI-map construction, and pseudo-label acceptance.
This makes LOGO a stricter test than simply training on non-test labelled
windows while leaving the held-out group's unlabelled pool in the augmentation
set.

\begin{table}[t]
\caption{Augmentation variants. All use the same RBF-SVM and target pipeline.}
\label{tab:variant_defs}
\centering
\setlength{\tabcolsep}{3pt}
\small
\begin{tabular}{l>{\raggedright\arraybackslash}p{0.60\columnwidth}}
\toprule
\textbf{Variant} & \textbf{Pseudo-label / weight rule} \\
\midrule
A0 & No augmentation; labelled T0+T1 windows only. \\
A1 & GP pseudo-labels for Valence and Dominance only. \\
A2 & GP pseudo-labels for all VAD dimensions. \\
A3 & EDA arousal labels + GP Valence/Dominance. \\
A5 & kNN arousal labels in 49-d feature space + GP V/D. \\
A6 & BFI-weighted GP V/D + BFI-weighted kNN arousal. \\
AP1 & GP class argmax, BFI cosine similarity as trust weight. \\
AP2 & GP class argmax, BFI similarity multiplied by GP confidence. \\
\bottomrule
\end{tabular}
\end{table}

\subsection{Why Personality Should Weight Trust}
\label{sec:why_bfi}

The AP1 design deliberately avoids using personality as a high-capacity neural
conditioning input.
At $N=103$, adding trait-conditioned projection matrices gives the model more
ways to overfit than the data can support.
Instead, personality plays a weaker role: it adjusts how much a pseudo-label
should be trusted.

This distinction matters conceptually.
Personality is not expected to predict the exact affective state of a
15-second physiological window.
It is expected to indicate whether two group members are likely to appraise the
same collaborative situation similarly.
Agreeableness, Extraversion, and Neuroticism can shape how participants
experience cooperation, conflict, and social pressure, but those effects are
slow and contextual.
Using BFI as a sample weight respects that timescale.

The hypothesis is therefore not ``personality predicts physiology.''
The hypothesis is: when pseudo-labels are propagated through a group, labels
coming from personality-similar peers are more trustworthy than labels coming
from dissimilar peers.
AP1 tests that hypothesis directly.

\paragraph{Capacity control.}
The SVM does not receive BFI features directly.
This prevents the model from memorising participant or group-level trait
patterns.
Personality only enters through sample weights on pseudo-labelled windows,
which constrains it to a calibration role.
This capacity control is one reason the AP1 result is interpretable despite
the small labelled set.

\section{Results}
\label{sec:results}

\subsection{Known-Team Augmentation}
\label{sec:known_team}

Table~\ref{tab:svm_aug} reports the central augmentation comparison.
The no-augmentation SVM reaches $0.495$ mean macro-F1.
GP augmentation improves modestly, especially when all dimensions are enabled.
EDA-derived arousal labels hurt performance, and lowering the confidence
threshold worsens the effect.
Personality-mediated augmentation performs best: AP1 reaches $0.526$, with
Valence and Dominance gaining most.

\begin{table}[t]
\caption{SVM augmentation variants on the canonical T3 test split.}
\label{tab:svm_aug}
\centering
\setlength{\tabcolsep}{4pt}
\small
\begin{tabular}{lccc}
\toprule
\textbf{Variant} & \textbf{V} & \textbf{A} & \textbf{D} \\
\midrule
A0: No aug              & 0.570 & 0.530 & 0.385 \\
A1: GP V/D only         & 0.592 & 0.530 & 0.406 \\
A2: GP all dims         & 0.592 & \best{0.546} & 0.406 \\
A3: EDA + GP V/D        & 0.458 & 0.530 & 0.430 \\
A5: kNN arousal + GP V/D& 0.592 & 0.521 & 0.406 \\
A6: BFI$\times$GP + kNN & 0.593 & 0.535 & 0.431 \\
\midrule
\best{AP1: BFI-only}    & \best{0.608} & 0.535 & \best{0.435} \\
AP2: BFI$\times$GP      & 0.593 & 0.544 & 0.431 \\
\bottomrule
\end{tabular}
\end{table}

The pattern supports a dimension-specific interpretation.
Valence and Dominance are appraisal-heavy dimensions whose labels are more
likely to transfer among personality-similar participants.
Arousal is phasic and event-driven; it benefits more from temporal proximity
to self-report than from personality similarity.
In mean macro-F1, AP1 improves over no augmentation by $+0.031$.
The gain is not large, but it is concentrated where the mechanism predicts:
Valence improves by $+0.038$ and Dominance by $+0.050$, while Arousal changes
little.
This makes AP1 a calibrated low-label improvement rather than a broad claim
that personality solves physiological affect recognition.

\subsection{Physiological Alternatives for Arousal}
\label{sec:arousal_alternatives}

Arousal is the hardest dimension for personality-weighted augmentation because
its short-timescale physiology is not primarily trait-like.
EDA and PPG respond to task events, movement, cognitive demand, and social
pressure, often within tens of seconds.
We therefore tested physiology-derived pseudo-labels as alternatives to GP
arousal.

The EDA-only regression path is interpretable but too compressed.
Ridge regression predicts near the mean of the SAM scale, which causes most
pseudo-labels to fall into the Mid class after binning.
Relaxing the confidence threshold accepts more windows but increases noise,
leading to the worst result in the augmentation table.
The kNN arousal path performs better because it uses the full 49-dimensional
feature space, including IMU, PPG, gaze, and pupil features.
However, kNN still fails to beat the GP arousal variant on the T3 test split.
This suggests that temporal proximity to self-report remains more reliable
than feature-space similarity for arousal at this label density.
A separate preprocessing analysis of the same benchmark reports that smoothing
EDA and PPG can stabilise arousal estimates, while preserving most valence
signal. In this paper, the augmentation finding is narrower: personality
weighting is most reliable for appraisal-heavy dimensions (Valence and
Dominance), while Arousal remains primarily constrained by temporal
proximity and signal quality.

\subsection{Personality Correlations as Mechanistic Support}
\label{sec:bfi_support}

Participant-level BFI correlations provide a useful sanity check for the AP1
mechanism.
Agreeableness correlates positively with both Valence and Arousal, while
Neuroticism correlates negatively with Arousal.
These correlations do not prove that BFI should be a direct model input; in
fact, direct BFI conditioning overfits.
They do support the weaker AP1 claim that personality similarity carries
information about affective appraisal in collaborative tasks.

\begin{table}[t]
\caption{Participant-level Spearman correlations between BFI traits and VAD.}
\label{tab:bfi_corr_short}
\centering
\setlength{\tabcolsep}{4pt}
\small
\begin{tabular}{lccc}
\toprule
\textbf{Trait} & \textbf{Valence} & \textbf{Arousal} & \textbf{Dominance} \\
\midrule
Extraversion      & $+0.075$ & $+0.062$ & $-0.065$ \\
Agreeableness     & $+0.416^{**}$ & $+0.423^{**}$ & $+0.269$ \\
Conscientiousness & $+0.234$ & $+0.304$ & $+0.108$ \\
Neuroticism       & $-0.049$ & $-0.349^{*}$ & $-0.058$ \\
Openness          & $+0.162$ & $+0.196$ & $+0.106$ \\
\bottomrule
\end{tabular}
\end{table}

\subsection{Strict Session Isolation}
\label{sec:logo}

Table~\ref{tab:logo} shows that the AP1 gain is scoped.
When the held-out group is excluded from all pool generation, augmentation no
longer improves performance.
This does not invalidate AP1; it identifies its deployment regime.
Personality weighting is useful for a known team with earlier-task data, not a
general cross-group transfer mechanism at 10-group scale.

The LOGO result is also a warning about small group datasets.
Adding pseudo-labelled windows from other groups increases data volume, but it
does not necessarily increase relevant information for a new group.
Groups develop idiosyncratic baselines, conversational rhythms, and shared
task interpretations.
At 10 groups, personality similarity across groups is too weak a bridge to
overcome those differences.

\begin{table}[t]
\caption{Session-isolated LOGO augmentation. Pool and BFI maps exclude the
held-out session.}
\label{tab:logo}
\centering
\setlength{\tabcolsep}{4pt}
\small
\begin{tabular}{lccc}
\toprule
\textbf{Variant} & \textbf{V} & \textbf{A} & \textbf{D} \\
\midrule
A0: No aug     & 0.235 & 0.307 & 0.254 \\
AP1: BFI-only  & 0.224 & 0.279 & 0.284 \\
A2: GP all dims& 0.218 & 0.291 & 0.278 \\
\bottomrule
\end{tabular}
\end{table}

\subsection{LOSO Cross-Subject Stress Test}
\label{sec:loso}

To separate cross-person from cross-group transfer, we additionally report a
LOSO stress test under the same balanced target-construction pipeline.
Each fold holds out one participant's T3 windows and aggregates results across
the 32 participants with T3 labels.
Following strict cross-subject isolation, no pseudo-label augmentation is used
in this stress test.
Aggregate macro-F1 is $0.320$ (V$=0.190$, A$=0.438$, D$=0.334$), confirming a
substantial transfer gap relative to known-team temporal holdout.

\begin{table}[t]
\caption{LOSO cross-subject stress test (no augmentation).
Aggregate over 32 held-out participants with T3 labels.}
\label{tab:loso}
\centering
\setlength{\tabcolsep}{4pt}
\small
\begin{tabular}{lccc}
\toprule
\textbf{Variant} & \textbf{V} & \textbf{A} & \textbf{D} \\
\midrule
A0: LOSO no aug & 0.190 & 0.438 & 0.334 \\
\bottomrule
\end{tabular}
\end{table}

\subsection{Uncertainty and Significance}
\label{sec:uncertainty}

The canonical T0+T1$\rightarrow$T3 comparison is a fixed split with
deterministic SVM training, so Table~\ref{tab:svm_aug} reports point
estimates rather than seed variance.
For protocol-level uncertainty, LOGO is evaluated across 10 held-out
sessions; mean macro-F1 is $0.266$ for A0, $0.262$ for AP1, and $0.262$ for
A2.
Paired Wilcoxon signed-rank tests on per-session mean macro-F1 show no
significant AP1/A2 difference against A0 (all $p>.2$), supporting the claim
that augmentation does not improve unseen-group transfer.

\subsection{External Questionnaire Validation}
\label{sec:questionnaire}

We aggregate SVM predictions to participant-by-task level and correlate them
with post-block questionnaire items not used in training.
Predicted Arousal correlates most strongly with T3 engagement ($r=0.46$),
while predicted Valence correlates with T3 satisfaction ($r=0.39$).
These correlations support the psychological validity of the target pipeline
and augmentation recipes.

The questionnaire results are important because the model is trained only on
VAD self-reports.
Engagement and satisfaction are independent constructs collected after the
task, so their correlations with predicted Arousal and Valence indicate that
the physiological predictions are not merely reproducing arbitrary label
boundaries.
The absence of strong Dominance correlations is also consistent with the
classification results: Dominance is the weakest VAD dimension and may require
relational features such as speaking turns, gaze-to-person, and posture.
Because four questionnaire associations are tested on small samples, these
results are exploratory and should be interpreted with uncorrected
multiple-comparison risk in mind.

\begin{table}[t]
\caption{Exploratory external validation against post-block questionnaire
items.}
\label{tab:q}
\centering
\setlength{\tabcolsep}{3pt}
\footnotesize
\begin{tabular}{@{}llllcc@{}}
\toprule
\textbf{Task} & \textbf{Pred.} & \textbf{Q item} & \textbf{$r$} & \textbf{$n$} \\
\midrule
T1 & Valence & overall valence & $+0.41$ & 33 \\
T3 & Valence & satisfaction    & $+0.39$ & 31 \\
T3 & Arousal & engagement      & $+0.46$ & 32 \\
T3 & Arousal & mental demand   & $+0.17$ & 31 \\
\bottomrule
\end{tabular}
\end{table}

\section{Mechanistic Interpretation}
\label{sec:mechanism}

The augmentation results are easiest to interpret by VAD dimension rather than
by mean F1 alone.
Personality similarity should help most when the affective label depends on
stable appraisal tendencies and least when the label reflects short-lived
physiological activation.
This is exactly the pattern observed in Table~\ref{tab:svm_aug}.

\paragraph{Valence.}
Valence improves most reliably under AP1.
In collaborative tasks, positive or negative appraisal depends not only on
physiological activation but also on how a participant interprets cooperation,
agreement, conflict, and task progress.
Those interpretations plausibly relate to trait-level social orientation.
AP1 does not claim that agreeable participants are always high-valence; rather,
it gives more weight to pseudo-labels from peers whose trait profile suggests
similar appraisal of the same group event.

\paragraph{Arousal.}
Arousal is the counterexample that keeps the claim honest.
The strongest AP1 gains are not in Arousal because the underlying signal is
more phasic and event-driven.
EDA and PPG can change quickly after cognitive demand, surprise, negotiation
pressure, movement, or public outcome reveal.
These changes are better captured by temporal proximity to self-report and by
preprocessing that respects autonomic timescales than by trait similarity.
This explains why GP-all-dimensions has the highest Arousal score in
Table~\ref{tab:svm_aug}, while AP1 is strongest on Valence and Dominance.

\paragraph{Dominance.}
Dominance is the most socially relational dimension.
It reflects perceived control, influence, and agency during group interaction.
AP1 improves Dominance because personality-similar peers may provide more
credible pseudo-labels for stable appraisal of control.
At the same time, Dominance remains hard: T4 Dominance is not administered,
and physiology-only features lack direct measures of who speaks, who is
addressed, who is looked at, and who controls the task outcome.
This motivates future integration of audio, gaze-to-person, posture, and turn
structure.

\paragraph{Why AP2 does not dominate AP1.}
AP2 multiplies BFI similarity by GP confidence, but this extra gate slightly
reduces the mean score relative to AP1.
The likely reason is that the confidence signal is too sparse in this
low-label setting.
GP confidence removes uncertain windows, but some uncertainty is useful:
pseudo-label augmentation needs enough diverse windows to regularise the SVM
decision boundary.
AP1 keeps the stable personality trust signal while admitting a broader pool.

\section{Recommended Use}
\label{sec:recommended}

Table~\ref{tab:recipes} summarises when each augmentation recipe should be
used.
The key point is that AP1 is not a universal default.
It is appropriate when earlier-task data from the same team is available and
the target is a later task from that team.
For unseen-group deployment, the recommended baseline is no augmentation or
strictly session-isolated augmentation.

\begin{table}[t]
\caption{Recommended label-augmentation recipes and intended scope.}
\label{tab:recipes}
\centering
\setlength{\tabcolsep}{3pt}
\small
\begin{tabular}{>{\raggedright\arraybackslash}p{0.16\columnwidth}>{\raggedright\arraybackslash}p{0.29\columnwidth}>{\raggedright\arraybackslash}p{0.45\columnwidth}}
\toprule
\textbf{Recipe} & \textbf{Use case} & \textbf{Interpretation} \\
\midrule
A0 & Conservative baseline & No pseudo-label assumptions \\
A2 & GP temporal labels & Affect continuity in known teams \\
AP1 & Known-team augmentation & BFI-weighted pseudo-label trust \\
AP2 & Conservative AP1 & BFI trust gated by GP confidence \\
LOGO A0 & New-group baseline & Strict unseen-group deployment \\
\bottomrule
\end{tabular}
\end{table}

\begin{table}[t]
\caption{Checks required before claiming leakage-aware augmentation.}
\label{tab:checks}
\centering
\setlength{\tabcolsep}{3pt}
\small
\begin{tabular}{>{\raggedright\arraybackslash}p{0.34\columnwidth}>{\raggedright\arraybackslash}p{0.54\columnwidth}}
\toprule
\textbf{Check} & \textbf{Why it matters} \\
\midrule
Train-fitted thresholds & Keeps hard and pseudo-labels in one class space. \\
Exclude test-task labels & Prevents direct temporal label leakage. \\
Regenerate LOGO pool & Removes held-out group from all augmentation signals. \\
Keep BFI out of SVM features & Preserves personality as trust, not identity. \\
Report V/A/D separately & Reveals that AP1 mainly helps V and D. \\
\bottomrule
\end{tabular}
\end{table}

\section{Discussion}
\label{sec:discussion}

\paragraph{Personality is better as a trust weight than as a direct input.}
Direct BFI conditioning overfits at the current label density, whereas AP1
uses BFI only to weight pseudo-label trust.
This is a lower-capacity role for personality and better matched to $N=103$.

\paragraph{Known-team benefit is still useful.}
Many practical group-sensing deployments observe the same team over multiple
tasks or sessions.
In that setting, AP1 is legitimate: it uses pre-session personality and
earlier-task unlabelled windows to smooth the team's affect trajectory.
The important wording is ``same team.''
The AP1 result should not be sold as generalisation to strangers; it is a
within-team calibration method for sparse labels.
That narrower claim is still useful because many collaborative systems care
about monitoring an ongoing team rather than classifying arbitrary new groups.

\paragraph{Unseen-group transfer remains unsolved.}
LOGO performance remains far below known-team performance.
This gap likely reflects group-specific baselines, interaction dynamics, and
role structures not captured by per-person physiology and BFI alone.
For new groups, a better strategy may be rapid calibration: collect one or two
short self-report anchors from the group, fit group-specific baselines, and
then apply AP1 only after the group has entered the known-team regime.
This hypothesis is not tested here, but it follows directly from the gap
between canonical temporal holdout and strict LOGO.

\paragraph{The SVM is a feature, not a weakness.}
The strongest results in this paper use an RBF-SVM rather than a neural model.
For $N=103$, this is appropriate.
The purpose is not to demonstrate architectural novelty; it is to isolate
whether a pseudo-label trust signal is useful when model capacity is held
constant.
The SVM makes the augmentation comparison more transparent.

\section{Limitations}
\label{sec:limitations}

The dataset contains only 10 groups, so strict LOGO estimates have high fold
variance.
The canonical AP1 comparison should be interpreted as a temporal known-team
result, not as unseen-group deployment.
Future work should add clustered confidence intervals, evaluate repeated
sessions per group, and test calibration strategies for new teams.
The reported AP1 improvement is small in absolute terms and should be treated
as a calibrated low-label gain rather than a large general-purpose improvement.
Because windows overlap within participants and groups, future statistical
tests should resample at participant or session level rather than assuming
independent windows.

\section{Conclusion}
\label{sec:conclusion}

Personality-weighted label augmentation improves physiological VAD prediction
when the deployment setting is a known team with sparse earlier-task labels.
BFI-44 cosine similarity is most useful as a pseudo-label trust weight,
especially for Valence and Dominance, but it does not solve unseen-group
transfer.
This boundary is the contribution: personality helps label augmentation under
the right social and evaluation regime.

\bibliographystyle{ACM-Reference-Format}
\bibliography{icmi2026_refs}

\end{document}
